\section{Discussion}

The experimental results demonstrate that architectural complexity correlates positively with multi-label recommendation quality up to a point, beyond which gains plateau. FusionNet's design choice of maintaining separate encoding pathways for symptom text and structured patient features proves particularly well-suited to the outpatient setting, where demographic and comorbidity signals carry meaningful clinical weight alongside free-text complaints. The explicit fusion of these two modalities, rather than simple concatenation at the embedding level as in MLP and CNN, enables the model to learn richer cross-modal interactions, reflected in its superior subset accuracy and Jaccard index.

The consistently high ROC-AUC scores across all models suggest that the underlying symptom--medication signal in the synthetic dataset is learnable even by shallow architectures. However, the large gap in subset accuracy between MLP (0.4210) and FusionNet (0.6455) underscores that correctly predicting the full medication set for a patient, rather than individual drug labels in isolation, is a substantially harder problem that benefits from sequential or attention-based inductive biases. This distinction is clinically significant: a system that recommends most drugs correctly but omits one critical medication or includes one unsafe drug is of limited utility in practice.

The low macro-averaged F1 scores observed uniformly across all models highlight a known challenge in medication recommendation: the long-tail distribution of drug classes. Infrequent medications, which may be the most critical in edge cases such as rare comorbidities or drug allergies, are systematically under-predicted. This limitation is partly an artifact of the synthetic data generation process, where drug assignment is driven by disease--drug mappings rather than physician behavior, and may underrepresent clinically nuanced prescribing decisions. Addressing this imbalance through class-weighted loss functions, oversampling strategies, or prototype-based learning for rare drug classes is an important direction for future work.

SelfAttention's strong PRAUC and Precision@K performance relative to its subset accuracy suggests that attention mechanisms are particularly effective at surfacing the most relevant drugs at the top of the ranked list, even when the full predicted set is not exact. This property is desirable in a clinical decision support context, where a ranked shortlist presented to the clinician is often more actionable than a fixed predicted set. Integrating self-attention as a ranking re-scorer on top of FusionNet's predictions is a promising architectural extension.

A key limitation of the current evaluation is the absence of an explicit DDI rate comparison across models at inference time, as the safety filtering layer described in the methodology was applied during dataset construction rather than as a post-hoc inference component in these experiments. Future work should decouple the safety filter and measure the raw DDI rate of model outputs before and after filtering, enabling a direct quantification of the safety benefit provided by the graph-based interaction layer. Additionally, all experiments were conducted on synthetically generated data, and performance on real-world clinical datasets may differ considerably due to the distributional gap between template-constructed symptom text and authentic clinical notes.

\section{Conclusion}

This paper presented a survey of twenty deep learning approaches to medication recommendation, identified persistent gaps in outpatient generalization, safety integration, and interpretability, and proposed a hybrid architecture addressing these limitations. Five neural models of increasing complexity were implemented and evaluated on a synthetic yet clinically grounded multi-label dataset constructed by integrating symptom--disease mappings, drug--disease associations, DDI constraints, and Synthea demographic records.

FusionNet, which combines a sequential text encoder with a dedicated structured feature branch, achieved the strongest overall performance across Jaccard index, micro-F1, and subset accuracy, demonstrating the value of modality-aware feature fusion for outpatient medication recommendation. SelfAttention showed competitive top-$k$ ranking performance, suggesting its suitability in clinician-facing shortlist generation. All models achieved high ROC-AUC scores, confirming that the symptom-to-medication mapping signal is robustly learnable, while low macro-F1 scores highlighted the difficulty of rare drug prediction as an open challenge.

The proposed framework integrates a DDI-aware safety filter and SHAP-based interpretability layer to address clinical trust and patient safety requirements beyond raw predictive performance. Together, these components form a foundation for a lightweight, safety-aware, and explainable medication recommendation system suitable for primary-care decision support under the data-sparse conditions typical of outpatient settings.

Future work will focus on evaluation on real-world clinical datasets, incorporation of class-imbalance mitigation strategies for rare drug classes, extension of the safety filter to runtime inference with measurable DDI rate reduction, and exploration of larger pretrained clinical language models as drop-in symptom encoders within the FusionNet backbone.
